% DRAFT 2026-11-28 (Tayyab's agent, from the project records). Only measured values; sources in brackets
% (R-NNN = docs/progress/TRAINING.md, D-NNN = DECISIONS.md, P&S = Problems_and_Solutions.md).
% TODO(agent, 2026-11-28 night): R-136 rows, layout fallback, viewer load times, .spz quality.
\chapter{Evaluation}
\label{ch:evaluation}

\section{Method}
\paragraph{Image quality.} Frames are held out of training (a fixed list per recording) and the trained model is
rendered at their camera poses. Because training refines the cameras, each held-out camera is first aligned to the
frozen model (six degrees of freedom plus a per-channel colour gain), then PSNR and SSIM are measured at half resolution
(\texttt{eval\_align.py}); the size of each correction is reported too. Older results use the colour-fitted PSNR
without pose alignment (\texttt{eval\_offpath.py}); the two are not mixed in one comparison.
\paragraph{Metric accuracy.} Tape-measured lengths against the same lengths in the model (\texttt{tape\_check.py}).
\paragraph{Time.} Every stage of every run writes its own duration (\texttt{timing.json}, \texttt{stages.tsv}); all runs
are listed in \texttt{docs/progress/runs.tsv} (D-090).
\paragraph{Machine.} Unless stated, an RTX~3070 with 8\,GB, 16\,GB of RAM, Windows~11.

\section{Phone Tracking against the Pipeline's Cameras}
\label{sec:eval-arcore}
The project's first change of plan (D-014) was to take the camera poses from SfM instead of ARCore. Table~\ref{tab:eval-arcore}
shows why, on the keyframes of the automatic runs: ARCore's own camera path lies 0.14 to 0.25\,m (median) from the
SfM cameras on textured rooms and tracks only a quarter of the featureless walk (R-141). The reference is the
pipeline's own cameras, not ground truth.
\begin{table}[htbp]
\centering\small
\caption{ARCore camera path against the pipeline's cameras (Sim3-aligned, position error)}
\label{tab:eval-arcore}
\begin{tabular}{p{3.4cm}p{1.6cm}p{2.2cm}p{1.8cm}p{3.4cm}}
\toprule
\textbf{Room} & \textbf{Keyframes} & \textbf{ARCore tracked} & \textbf{Placed} & \textbf{Error median / p90 (m)} \\
\midrule
Main test room & 1005 & 100\,\% & 1001 & 0.245 / 0.690 \\
Bedroom & 1060 & 88\,\% & 1025 & 0.138 / 0.337 \\
scan\_009 (featureless) & 894 & 26\,\% & 727$^g$ & 1.33 / 3.68 \\
\bottomrule
\end{tabular}

\smallskip
\footnotesize $^g$ placed by the sensor-assisted chain; its units are metric only approximately.
\end{table}
The second half of the exit-gate comparison is SfM with and without the ARCore-guided image pairs, on the same
keyframes (R-142). Both versions placed the same frames: all 1005 on the main room, 409 of 894 on scan\_009, in every
15-second stretch of the walk. On the main room the two solutions differ by 0.23\,m (median) along the walk. With image
retrieval in place, ARCore's contribution is therefore the metric scale and gravity of the model, not more frames
tracked.

\section{Reconstruction Quality}
\begin{table}[htbp]
\centering\small
\caption{Held-out image quality on the main test room (newRoomAttemptSadness, 240-second video)}
\label{tab:eval-quality}
\begin{tabular}{p{4.6cm}p{2.2cm}p{1.6cm}p{1.4cm}p{3.2cm}}
\toprule
\textbf{Model} & \textbf{Held-out frames} & \textbf{PSNR (dB)} & \textbf{SSIM} & \textbf{Source} \\
\midrule
v6 (mid report) & 42 & 26.52 & 0.8886 & mid report \\
v9 end to end & 79 & 25.16$^a$ & 0.8247 & P\&S 4.16 \\
v9, faster training & 79 & 25.14$^a$ & 0.8251 & R-062 \\
v9, faster training + photometric init & 79 & 25.16$^a$ & 0.8240 & R-063 \\
Automatic pipeline (\texttt{--auto}) & 62$^b$ & 24.61$^a$ (25.33 aligned) & 0.8222 (0.8350) & R-136 \\
Automatic, photometric initialisation & 62$^b$ & 24.50$^a$ (25.11 aligned) & 0.8233 (0.8335) & R-144 \\
\bottomrule
\end{tabular}

\smallskip
\footnotesize $^a$ colour-fitted PSNR without pose alignment. $^b$ the v9 list mapped to the automatic run's keyframes
(49 of 79 identical, the rest at most 0.2\,s away).
\end{table}
The v6 and v9 rows use different held-out lists and are not directly comparable: the v9 list holds out off-path
views, which are harder.

\section{Featureless Walls}
\begin{table}[htbp]
\centering\small
\caption{scan\_009 (featureless walls): aligned PSNR on the same 63 held-out frames}
\label{tab:eval-featureless}
\begin{tabular}{p{6.4cm}p{1.8cm}p{1.6cm}p{2.0cm}}
\toprule
\textbf{Cameras and training} & \textbf{PSNR (dB)} & \textbf{SSIM} & \textbf{Source} \\
\midrule
Normal SfM & \multicolumn{3}{p{5.4cm}}{87 of 5{,}629 frames placed; no usable model (R-045)} \\
Sensor chain with scale-drift correction & 22.67 & & R-094 / R-096 \\
+ loop closure with long-range pairs & 23.63 & 0.888 & R-096 \\
+ camera refinement in training & 23.96 & 0.896 & R-120 \\
+ extra video frames in weak stretches & 24.25 & 0.898 & R-128 \\
Same, Depth Anything 3 priors with confidence cut-off$^c$ & 24.25 & & R-133 \\
Automatic pipeline (\texttt{--auto}, sensor path chosen by itself) & 23.07$^d$ & & R-135 \\
Same + track directions in the rotation solve (rejected) & 21.81$^h$ & & R-140 \\
\textbf{Automatic pipeline with ALIKED features} & \textbf{25.68}$^d$ & & R-145 \\
\bottomrule
\end{tabular}

\smallskip
\footnotesize $^c$ commercial-use weights (D-098); 26.56\,dB on the blank-wall stretch (60--90\,s), the best of all runs.
$^d$ on the 61 frames it shares with R-128 (R-128: 24.46 on the same 61). $^h$ on the 56 frames shared by R-140,
R-135, R-128 and R-133 (23.03, 24.33 and 24.41 on the same 56).
\end{table}
The first automatic run reached the sensor-assisted chain without a hand choice (SfM placed 41\,\% of the keyframes;
after the chain 81\,\%), but its cameras were worse than the hand-built chain's, by 1.4\,dB overall and 3\,dB on the
blank wall (R-135). Its camera positions drifted apart from the hand-built ones from 45\,s on, and its loop closure had
half the image-pair directions: the hand-built chain had been matched with learned ALIKED features, the automatic one
with SIFT. With ALIKED matching the automatic pipeline placed 78\,\% of the keyframes in SfM alone and became the best
scan\_009 model of the project: 25.68\,dB against 24.46\,dB for the best hand-built model on the same 61 frames,
28.66\,dB on the blank wall and 22.05\,dB on the 134--140\,s stretch that every earlier run failed (R-145). SfM with
ALIKED is slower (20.5 against 3.9 minutes). The room layout of this recording is still not trustworthy
(Section~\ref{sec:eval-layout}).

\section{The Automatic Pipeline}
\begin{table}[htbp]
\centering\small
\caption{Automatic pipeline (\texttt{run\_full.sh -- --auto}): raw video to web bundle, no hand-set option}
\label{tab:eval-auto}
\begin{tabular}{p{2.6cm}p{3.2cm}p{3.6cm}p{4.2cm}}
\toprule
\textbf{Recording} & \textbf{Analyzer choice} & \textbf{Time (min)} & \textbf{Result} \\
\midrule
Bedroom (textured), R-134 & normal SfM (97\,\% placed); 2 of 59 windows densified & model 83.5, layout 5.0, 2D labels 12.7, rest 2.8; 104.0 in all & 1.43\,M Gaussians, 5 walls, 19 fixtures, mesh 0.89\,M triangles, 33\,MB bundle; outline not closed \\
scan\_009 (featureless), R-135 & sensor path (41\,\% placed); 27 of 38 windows densified & model 56.4 (sensor chain 17.6), after 15.4; 71.9 in all & 23.07\,dB vs 24.46 (hand chain); layout failed (room height 7.28\,m) \\
Main test room, R-136 & normal SfM (100\,\% placed); 5 of 48 windows densified & model 86.0, after 20.1; 106.1 in all & 25.33\,dB aligned; closed room 6.575 $\times$ 4.380\,m, tape $-1.86$\,\%; 17 fixtures (unchecked) \\
\bottomrule
\end{tabular}
\end{table}

\section{Processing Time}
\begin{table}[htbp]
\centering\small
\caption{Pipeline time per stage (minutes)}
\label{tab:eval-time}
\begin{tabular}{p{3.0cm}p{1.2cm}p{1.0cm}p{1.0cm}p{1.0cm}p{1.0cm}p{1.2cm}p{1.2cm}p{1.2cm}}
\toprule
\textbf{Run} & \textbf{video (s)} & \textbf{SfM} & \textbf{priors} & \textbf{init} & \textbf{train} & \textbf{model} & \textbf{after} & \textbf{total} \\
\midrule
v9, main room (3070) & 240 & 22.4 & 8.1 & 20.4 & 42.5 & 96 & & \\
R-031, classroom (RTX 4000 Ada) & 185 & & & & & 64.7 & & 81.0 \\
R-134, bedroom, \texttt{--auto} & & 18.2 & 2.3 & 19.3 & 36.6 & 83.5 & 20.5 & 104.0$^f$ \\
R-135, scan\_009, \texttt{--auto} & 188 & 3.9 + 17.6$^e$ & 1.8 & 1.4 & 27.2 & 56.4 & 15.4 & 71.9 \\
R-136, main room, \texttt{--auto} & 240 & 28.4 & 1.9 & 17.5 & 32.9 & 86.0 & 20.1 & 106.1 \\
\bottomrule
\end{tabular}

\smallskip
\footnotesize $^e$ SfM plus the sensor-assisted camera chain. $^f$ sum of the stage times; the run was paused
twice. Empty cells: not recorded for that run.
\end{table}
NFR3 (a 60-second video in 15 minutes on an RTX~4090) is not met and was not measured on a 4090. A known saving not yet
in the automatic recipe: the photometric initialisation, 17.5 $\rightarrow$ 7.5\,min inside the automatic pipeline
(7.1\,min net, as training took 2.9\,min longer) at $-0.22$\,dB on the same held-out frames (R-144; R-063 on the older
recipe: 20.1 $\rightarrow$ 7.5\,min at equal quality); open decision OD-10.

\section{Metric Accuracy}
\label{sec:eval-accuracy}
On the v3 model, six tape-measured lengths from 0.46 to 6.70\,m differed by 4.8\,mm on average (0.19\,\%) and at most
0.29\,\%, well inside NFR6 ($\leq$\,5\,\%). The faster v9 SfM made the main room 6.617\,m long (v9 end to end; 6.60\,m in the all-frames model) against
6.70\,m on the tape ($-1.24$ to $-1.5$\,\%), still inside NFR6; its scale changes between SfM runs (known limit).
With the automatic pipeline (R-136) the main room comes out as a closed room of 6.575 $\times$ 4.380\,m; the tape
gives 6.700\,m for its length ($-1.86$\,\%, inside NFR6; the hand-set v9: 6.617\,m, $-1.24$\,\%).
% TODO(Umair): tape sheets for scans #3, #6, #7, #8 (issue #20).

\section{Room Layout}
\label{sec:eval-layout}
The layout found the walls, corners and height of the textured rooms: the main room as a closed four- to five-wall
room (2.70 to 2.83\,m high depending on the code version; 2.81\,m in R-136), the bedroom (2.78 to 3.02\,m; 2.94\,m in
R-134), the classroom (2.52\,m, R-035). It fails when the floor is not reconstructed: in scan\_009 a layer of floaters
3.4\,m below the cameras was taken as the floor (7.28\,m room, R-135). A plausibility check (camera 0.5--2.2\,m above
the floor, room 2--5\,m high) now replaces such a floor and leaves the textured rooms unchanged (R-137), but scan\_009
still comes out 5.25\,m high: floaters make its walls look 4.4\,m tall and the rule that the ceiling may not lie below
the wall tops (D-072) raises it again. Its height is reported with a warning and is not trustworthy.

\section{Semantic Labels}
On the main test room, 15 of 15 detected fixtures were correct and 15 of 17 visible fixtures were found (v10, P\&S
4.18); the rules were tuned on that room, so this is not a general accuracy figure. On other rooms the pipeline found
21 (classroom, R-031), 11 (scan~\#6, R-059), 22 (bedroom, R-069) and 19 fixtures (bedroom, automatic, R-134); these have
not yet been checked by a person, so precision and recall on unseen rooms (E3) are not measured.
% TODO(Umair): fixture check sheets (issue #21) -> precision / recall per room.

\section{Web Viewer and Compression}
\paragraph{Frame times.} On scan~\#7 (brick walls), dragging for six seconds gave 168 frames
with a worst frame of 195\,ms before the fix and 1{,}174 frames with a worst frame of 15\,ms after it (desktop Chrome on
the RTX~3070 PC; scan~\#8: 290 $\rightarrow$ 20\,ms). NFR5 ($\geq$\,30\,FPS in a mid-range phone browser) is not
measured on a phone.
With the fix, the best model (0.6 million Gaussians) loaded in 2.2\,s from the local server and scans \#7 and \#8 in 6.1
and 7.6\,s from the online bucket; idle, dragging and hovering ran at the headless browser's 5\,ms frame pacing with
no frame above 15\,ms and no long main-thread task (R-143, desktop).
% TODO(Zaid): phone measurement (E5).
\paragraph{Compression.} The web bundle stores every Gaussian of the model as \texttt{.spz} (quantised positions,
opacity, colour, scale, rotation and third-degree spherical harmonics, gzip). Decoding the bundle and rendering it from
the same cameras as the uncompressed model gives the loss from compression alone (Table~\ref{tab:eval-spz}, R-139):
the files are 92.3 to 92.8\,\% smaller (NFR4: $\geq$\,60\,\%) and their renders differ from the original's by
35.8 to 36.5\,dB PSNR on average (SSIM 0.98 to 0.99), which is not visible.
\begin{table}[htbp]
\centering\small
\caption{Web \texttt{.spz} against the trained \texttt{.ply}: size and render difference (R-139)}
\label{tab:eval-spz}
\begin{tabular}{p{3.6cm}p{1.6cm}p{1.2cm}p{1.2cm}p{1.4cm}p{2.4cm}p{1.2cm}}
\toprule
\textbf{Model} & \textbf{Gaussians} & \textbf{.ply MB} & \textbf{.spz MB} & \textbf{smaller} & \textbf{PSNR mean / min} & \textbf{SSIM} \\
\midrule
Bedroom, automatic (R-134) & 1{,}425{,}538 & 353.5 & 25.6 & 92.8\,\% & 36.12 / 30.90 & 0.990 \\
Bedroom, hand-set (R-069) & 1{,}445{,}529 & 358.5 & 26.0 & 92.7\,\% & 36.52 / 30.42 & 0.990 \\
Main room, best model & 605{,}481 & 150.2 & 11.5 & 92.3\,\% & 35.82 / 34.11 & 0.983 \\
Main room, automatic (R-136) & 637{,}000 & 158.0 & 12.1 & 92.3\,\% & 36.49 / 34.48 & 0.983 \\
\bottomrule
\end{tabular}
\end{table}
% TODO(Zaid): .ksplat numbers on the same models (OD-1).

\section{Capture}
NFR1 (30\,FPS) is not met: 28.9\,FPS on average with 3.3\,\% dropped-frame gaps on a 185-second recording (R-031).
NFR2 (quality check under 16\,ms per frame) is logged by the app but not yet measured on a phone.

\section{Requirements Summary}
\begin{table}[htbp]
\centering\small
\caption{Non-functional requirements}
\label{tab:eval-nfr}
\begin{tabular}{p{1.4cm}p{5.0cm}p{3.6cm}p{2.0cm}}
\toprule
\textbf{Req.} & \textbf{Target} & \textbf{Measured} & \textbf{Met} \\
\midrule
NFR1 & capture at $\geq$\,30\,FPS & 28.9\,FPS, 3.3\,\% gaps & no \\
NFR2 & quality check $<$\,16\,ms per frame & not measured & -- \\
NFR3 & 60\,s video in 15\,min on an RTX~4090 & 71.9--106.1\,min for 188--240\,s videos on an RTX~3070 & no \\
NFR4 & compression $\geq$\,60\,\% & 92\,\% (\texttt{.spz}) & yes \\
NFR5 & $\geq$\,30\,FPS in a mid-range browser & desktop only (worst frame 15\,ms) & not on a phone \\
NFR6 & measurement error $\leq$\,5\,\% & 0.19\,\% mean (v3); room length $-1.24$\,\% (v9), $-1.86$\,\% (automatic) & yes \\
\bottomrule
\end{tabular}
\end{table}
